% 4.7 诱导表示的直积
\section{诱导表示的直积}

在前几节中我们已经看到：若 $\mathbf{G}$ 是群而 $\mathbf{T}$ 是 $\mathbf{G}$ 的不变 Abel 子群，则 $\mathbf{G}$ 的表示可以由 $\mathbf{T}$ 的表示确定；具体做法已在 4.3 节开头概述。本节将说明如何确定 $\mathbf{G}$ 的两个表示的内直积。也就是说，若 $(\Gamma^{i}_{p} \uparrow \mathbf{G})$ 与 $(\Gamma^{j}_{q} \uparrow \mathbf{G})$ 是 $\mathbf{G}$ 的任意两个表示，我们将说明如何确定 Clebsch--Gordan 分解中的系数 $C^{ij, l}_{pq, r}$：
\begin{equation*}
(\Gamma^{i}_{p} \uparrow \mathbf{G}) \boxtimes (\Gamma^{j}_{q} \uparrow \mathbf{G}) \equiv \sum_{l}\sum_{r}C^{ij, l}_{pq, r}(\Gamma^{l}_{r} \uparrow \mathbf{G}).
\tag{4.7.1}
\end{equation*}
为了推进这一问题，有必要重新表述 1.3 节与 4.1 节的若干结果，并引入一些新概念，特别是\textit{双陪集}（double coset）。

\noindent\textbf{定义 4.7.1}\quad （复共轭表示）\ 设 $\mathbf{G}$ 是群，$\Gamma$ 是 $\mathbf{G}$ 的幺正表示。则由
\begin{equation*}
\Gamma^{*}(g) = \{\Gamma(g)\}^{*}
\tag{4.7.2}
\end{equation*}
定义的表示 $\Gamma^{*}$ 称为 $\Gamma$ 的\textit{复共轭表示}。

\noindent\textbf{定理 4.7.1}\quad 设 $\mathbf{G}$ 是群，其表示记为 $\Gamma^{i}$，并设 $\Gamma$ 等价于直和 $\sum^{r}_{i=1}c_{i}\Gamma^{i}$，则 $c_{i}$ 是 $\mathbf{A}(\mathbf{G})$（$\mathbf{G}$ 的全对称表示）在 $\Gamma^{*} \boxtimes \Gamma^{i}$（或 $\Gamma \boxtimes \Gamma^{i*}$）中出现的频率。

这几乎是显然的。由式 (1.3.18)，
\begin{equation*}
c_{i} = \frac{1}{|\mathbf{G}|}\sum_{g}\chi^{*}(g)\chi^{i}(g)
\tag{4.7.3}
\end{equation*}
\begin{equation*}
= \frac{1}{|\mathbf{G}|}\sum_{g}\chi^{\Gamma^{*} \boxtimes \Gamma^{i}}(g)
\tag{4.7.4}
\end{equation*}
（表示直积的特征标见式 (1.3.21)）。但 $\chi^{A}(g) = 1$ 对一切 $g \in \mathbf{G}$ 成立，故式 (4.7.4) 成为
\begin{equation*}
c_{i} = \frac{1}{|\mathbf{G}|}\sum_{g}\chi^{A*}(g)\chi^{\Gamma^{*} \boxtimes \Gamma^{i}}(g),
\tag{4.7.5}
\end{equation*}
反用式 (1.3.18) 即知它是 $\mathbf{A}(\mathbf{G})$ 在 $\Gamma^{*} \boxtimes \Gamma^{i}$ 中的频率。

由定理 4.7.1 可知，$C^{ij, l}_{pq, r}$ 是三重直积
\begin{center}
$\left(\Gamma^{i}_{p} \uparrow \mathbf{G}\right) \boxtimes \left(\Gamma^{j}_{q} \uparrow \mathbf{G}\right) \boxtimes \left(\Gamma^{l*}_{r} \uparrow \mathbf{G}\right)$
\end{center}
中全对称表示的频率。因此我们将研究这个等价的、而且事实证明更容易处理的问题。

下面重新表述 4.1 节的一个结果。设 $\mathbf{G}$ 是群，$\mathbf{K}$ 是任一具有表示 $\Lambda$ 的子群，则诱导表示 $\Lambda \uparrow \mathbf{G}$ 由式 (4.1.4) 定义：
\begin{equation*}
g\phi_{\tau r} = \sum_{t}\phi_{\gamma t}\Lambda(k_{m})_{tr}.
\tag{4.7.6}
\end{equation*}
这里 $g$ 是左陪集分解 $\mathbf{G} = \sum_{\sigma}p_{\sigma}\mathbf{K}$ 中的任一元素 $p_{\lambda}k_{s}$，且 $k_{m} = p^{-1}_{\gamma}gp_{\tau}$。由 $gp_{\tau} \in p_{\gamma}\mathbf{K}$ 的事实定义 $p_{g\tau} = p_{\gamma}$，则
\begin{equation*}
(\Lambda \uparrow \mathbf{G})(g)_{\gamma\tau} = \Lambda(p^{-1}_{\gamma}gp_{\tau})\delta_{\gamma, (g\tau)}.
\tag{4.7.7}
\end{equation*}
设 $\Lambda$ 的特征标为 $\psi$，$\Lambda \uparrow \mathbf{G}$ 的特征标为 $\chi$，则
\begin{equation*}
\chi(g) = \sum_{\sigma}\mathrm{Tr}\left\{(\Lambda \uparrow \mathbf{G})(g)_{\sigma\sigma}\right\}
\tag{4.7.8}
\end{equation*}
\begin{equation*}
= \sum_{\sigma}\mathrm{Tr}\left\{\Lambda(p^{-1}_{\sigma}gp_{\sigma})\right\}\delta_{\sigma, (g\sigma)}
\tag{4.7.9}
\end{equation*}
\begin{equation*}
= \sum_{\sigma}\psi(p^{-1}_{\sigma}gp_{\sigma})\delta_{\sigma, (g\sigma)}.
\tag{4.7.10}
\end{equation*}
注意对 $\sigma$ 的求和只取那些使 $p^{-1}_{\sigma}gp_{\sigma} \in \mathbf{K}$ 的陪集，即那些使 $p_{\sigma}\mathbf{K}p^{-1}_{\sigma} = \mathbf{K}_{\sigma}$ 含有 $g$ 的 $\sigma$。

顺带地，公式 (4.7.10) 证明了式 (4.1.1) 后评注的脚注的合理性。假设我们选取了另一组陪集代表 $q_{\sigma} = p_{\sigma}k_{\sigma}$（$k_{\sigma} \in \mathbf{K}$），那么特征标公式将成为
\begin{equation*}
\sum_{\sigma}\psi(q^{-1}_{\sigma}gq_{\sigma})\delta_{\sigma, (g\sigma)}.
\tag{4.7.11}
\end{equation*}
但 $q_{\sigma}\mathbf{K}q^{-1}_{\sigma} = p_{\sigma}k_{\sigma}\mathbf{K}k^{-1}_{\sigma}p^{-1}_{\sigma} = p_{\sigma}\mathbf{K}p^{-1}_{\sigma} = \mathbf{K}_{\sigma}$，且 $q^{-1}_{\sigma}gq_{\sigma}$ 与 $p^{-1}_{\sigma}gp_{\sigma}$ 属于 $\mathbf{K}$ 的同一元素（二者相差 $\mathbf{K}$ 中元素的共轭），于是两个公式给出相同的值。

注意公式 (4.7.10) 不论 $\Lambda$ 在 $\mathbf{K}$ 下是否可约都成立。这一点很重要，因为我们不仅关心 $\mathbf{K}$ 是小群、$\Lambda$ 是小表示（因而不可约）的情形，也关心 $\Lambda$ 一般可约、$\mathbf{K}$ 未必是小群的情形。

\noindent\textbf{定理 4.7.2}\quad （频率定理）\ 采用上述记号，$\mathbf{A}(\mathbf{G})$ 在 $\Lambda \uparrow \mathbf{G}$ 中的频率等于 $\mathbf{A}(\mathbf{K})$ 在 $\Lambda$ 中的频率。

记 $f(\Gamma^{i}) \mid \Gamma$ 为 $\Gamma^{i}$ 在 $\Gamma$ 中的频率。由式 (4.7.3) 与 (4.7.10)，
\begin{equation*}
f\left\{A(\mathbf{G})\right\} \mid (\Lambda \uparrow \mathbf{G}) = \frac{1}{|\mathbf{G}|}\sum_{g \in \mathbf{G}}\sum_{\sigma}\psi(p^{-1}_{\sigma}gp_{\sigma})\delta_{\sigma, (g\sigma)}.
\tag{4.7.12}
\end{equation*}
交换求和次序，这成为
\begin{equation*}
f\left\{A(\mathbf{G})\right\} \mid (\Lambda \uparrow \mathbf{G}) = \frac{1}{|\mathbf{G}|}\sum_{\sigma}\sum_{g \in \mathbf{K}_{\sigma}}\psi(p^{-1}_{\sigma}gp_{\sigma})
\tag{4.7.13}
\end{equation*}
\begin{equation*}
= \frac{1}{|\mathbf{G}|}\sum_{\sigma}\sum_{k \in \mathbf{K}}\psi(k).
\tag{4.7.14}
\end{equation*}
共有 $|\mathbf{G}|/|\mathbf{K}|$ 个陪集 $\sigma$，于是
\begin{equation*}
= \frac{1}{|\mathbf{K}|}\sum_{k \in \mathbf{K}}\psi(k)
\tag{4.7.15}
\end{equation*}
\begin{equation*}
= f\left\{A(\mathbf{K})\right\} \mid \Lambda,
\tag{4.7.16}
\end{equation*}
定理得证。

\noindent\textbf{定理 4.7.3}\quad 设 $\Lambda$ 是 $\mathbf{K}$ 的表示（特征标 $\psi$），$\Gamma$ 是 $\mathbf{G}$ 的表示（特征标 $\chi$），其中 $\mathbf{K}$ 是 $\mathbf{G}$ 的子群，则
\begin{equation*}
\Gamma \boxtimes (\Lambda \uparrow \mathbf{G}) \equiv \left[(\Gamma \downarrow \mathbf{K}) \boxtimes \Lambda\right] \uparrow \mathbf{G}.
\tag{4.7.17}
\end{equation*}
我们通过证明两个表示的特征标相同来建立这一等价。$g$ 在 $\Gamma \boxtimes (\Lambda \uparrow \mathbf{G})$ 中的特征标是
\begin{equation*}
\chi(g)\sum_{\sigma}\psi(p^{-1}_{\sigma}gp_{\sigma})\delta_{\sigma, (g\sigma)},
\end{equation*}
而 $g$ 在 $\left[(\Gamma \downarrow \mathbf{K}) \boxtimes \Lambda\right] \uparrow \mathbf{G}$ 中的特征标是
\begin{equation*}
\sum_{\sigma}\chi(p^{-1}_{\sigma}gp_{\sigma})\psi(p^{-1}_{\sigma}gp_{\sigma})\delta_{\sigma, (g\sigma)}.
\end{equation*}
但 $g$ 与 $p^{-1}_{\sigma}gp_{\sigma}$ 在 $\mathbf{G}$ 中同类，故对一切 $\sigma$ 有 $\chi(g) = \chi(p^{-1}_{\sigma}gp_{\sigma})$。结论得证。

现在证明诱导的传递性定理，比引理 4.3.3 更一般的情形。

\noindent\textbf{定理 4.7.4}\quad （诱导的传递性）\ 设 $\mathbf{K}$ 是 $\mathbf{H}$ 的子群，$\mathbf{H}$ 是 $\mathbf{G}$ 的子群，$\Lambda$ 是 $\mathbf{K}$ 的表示。则
\begin{equation*}
(\Lambda \uparrow \mathbf{H}) \uparrow \mathbf{G} \equiv \Lambda \uparrow \mathbf{G}.
\tag{4.7.18}
\end{equation*}
我们通过证明：若 $\Gamma$ 是 $\mathbf{G}$ 的表示，则它在 $(\Lambda \uparrow \mathbf{H}) \uparrow \mathbf{G}$ 中的频率与它在 $\Lambda \uparrow \mathbf{G}$ 中的频率相同，来建立等价。
\begin{equation*}
\begin{array}{lll}
f(\Gamma) \mid (\Lambda \uparrow \mathbf{G}) & &\\
= f\left\{A(\mathbf{G})\right\} \mid \Gamma^{*} \boxtimes (\Lambda \uparrow \mathbf{G}), & & \text{由定理 4.7.1}\\
= f\left\{A(\mathbf{G})\right\} \mid \left[(\Gamma^{*} \downarrow \mathbf{K}) \boxtimes \Lambda\right] \uparrow \mathbf{G}, & & \text{由定理 4.7.3}
\end{array}
\end{equation*}
但
\begin{equation*}
\begin{array}{lll}
f(\Gamma) \mid (\Lambda \uparrow \mathbf{H}) \uparrow \mathbf{G} & &\\
= f\left\{A(\mathbf{G})\right\} \mid \Gamma^{*} \boxtimes \left[(\Lambda \uparrow \mathbf{H}) \uparrow \mathbf{G}\right], & & \text{由定理 4.7.1}\\
= f\left\{A(\mathbf{G})\right\} \mid \left[(\Gamma^{*} \downarrow \mathbf{H}) \boxtimes \Lambda\right] \uparrow \mathbf{G}, & & \text{由定理 4.7.3}\\
= f\left\{A(\mathbf{G})\right\} \mid \left[(\Gamma^{*} \downarrow \mathbf{K}) \boxtimes \Lambda\right] \uparrow \mathbf{G}, & & \text{由限制的传递性}
\end{array}
\end{equation*}
且由于限制的传递性 $(\Gamma^{*} \downarrow \mathbf{H}) \downarrow \mathbf{K} \equiv (\Gamma^{*} \downarrow \mathbf{K})$——这是显然的，因为两端由同样的矩阵组成——定理得证。

\noindent\textbf{定义 4.7.2}\quad 设 $\mathbf{H}$ 与 $\mathbf{K}$ 是 $\mathbf{G}$ 的任意两个子群。则
\begin{equation*}
\mathbf{G} = \sum_{\alpha}\mathbf{H}d_{\alpha}\mathbf{K},
\tag{4.7.19}
\end{equation*}
（在复形 $\mathbf{H}d_{\alpha}\mathbf{K}$ 中，每个元素只计一次，无论它出现多少次）称为 $\mathbf{G}$ 关于 $\mathbf{H}$ 与 $\mathbf{K}$ 的\textit{双陪集分解}。$d_{\alpha}$ 称为双陪集代表。

展开 (4.7.19) 是唯一的，方式与普通的单陪集分解大致相同。双陪集代表不唯一：双陪集中的任何元素都可以同样好地充当其代表。设 $g \in \mathbf{H}d_{\alpha}\mathbf{K}$，则 $g = hd_{\alpha}k$（$h \in \mathbf{H}$，$k \in \mathbf{K}$），且该陪集由所有这种形式的元素组成。

\noindent\textbf{定理 4.7.5}\quad 双陪集 $\mathbf{H}d_{\alpha}\mathbf{K}$ 含有 $|\mathbf{H}|/|\mathbf{L}_{\alpha}|$ 个 $\mathbf{K}$ 的左陪集（从而含有 $|\mathbf{H}||\mathbf{K}|/|\mathbf{L}_{\alpha}|$ 个元素），其中 $\mathbf{L}_{\alpha} = \mathbf{H} \cap d^{-1}_{\alpha}\mathbf{K}d_{\alpha}$。

证明如下。设
\begin{equation*}
\mathbf{H} = \sum_{\beta}q_{\alpha\beta}\mathbf{L}_{\alpha}.
\tag{4.7.20}
\end{equation*}
设 $h_{1}$ 与 $h_{2}$ 属于 $\mathbf{L}_{\alpha}$，则对任意 $\beta$ 有 $q_{\alpha\beta}h_{1}d_{\alpha}\mathbf{K} = q_{\alpha\beta}h_{2}d_{\alpha}\mathbf{K}$。这是因为可以写 $h_{1} = d_{\alpha}k_{1}d^{-1}_{\alpha}$、$h_{2} = d_{\alpha}k_{2}d^{-1}_{\alpha}$（$k_{1}, k_{2} \in \mathbf{K}$），且 $q_{\alpha\beta}h_{2}^{-1}h_{1}q_{\alpha\beta} \in \mathbf{K}$。因此形如 $q_{\alpha\beta}hd_{\alpha}\mathbf{K}$（$\beta$ 取遍式 (4.7.20)）的左陪集两两不同，且它们的数目是 $|\mathbf{H}|/|\mathbf{L}_{\alpha}|$；另一方面，$\mathbf{H}d_{\alpha}\mathbf{K}$ 中每个元素都属于某个这样的陪集。定理得证。

\noindent\textbf{定理 4.7.6}\quad （Mackey 子群定理）
\begin{equation*}
(\mathbf{D} \uparrow \mathbf{G}) \downarrow \mathbf{H} \equiv \sum_{\alpha}(\mathbf{D}_{\alpha} \downarrow \mathbf{L}_{\alpha}) \uparrow \mathbf{H}.
\tag{4.7.21}
\end{equation*}
回忆分解 $\mathbf{G} = \sum_{\sigma}p_{\sigma}\mathbf{K}$、$\mathbf{G} = \sum_{\alpha}\mathbf{H}d_{\alpha}\mathbf{K}$ 与 $\mathbf{H} = \sum_{\beta}q_{\alpha\beta}\mathbf{L}_{\alpha}$，并注意
\begin{equation*}
\chi(h) = \sum_{\sigma}\psi(p^{-1}_{\sigma}hp_{\sigma})\delta_{\sigma, (h\sigma)}.
\tag{4.7.22}
\end{equation*}
即对 $\sigma$ 的求和限制在使 $h \in \mathbf{K}_{\sigma}$ 的那些 $\sigma$ 上。现在 $\mathbf{K}_{\sigma} = p_{\sigma}\mathbf{K}p^{-1}_{\sigma}$，而每个 $p_{\sigma}$ 可以写成 $q_{\alpha\beta}d_{\alpha}k$ 的形式，于是
\begin{equation*}
\chi(h) = \sum_{\alpha, \beta}\psi\left(d^{-1}_{\alpha}q^{-1}_{\alpha\beta}hq_{\alpha\beta}d_{\alpha}\right)
\tag{4.7.23}
\end{equation*}
其中对 $\alpha, \beta$ 的求和限制在使 $q_{\alpha\beta}\mathbf{K}_{\alpha}q^{-1}_{\alpha\beta} \supset h$ 的那些 $\alpha, \beta$ 上。

现在
\begin{equation*}
\mathbf{D}(d^{-1}_{\alpha}q^{-1}_{\alpha\beta}hq_{\alpha\beta}d_{\alpha}) = \mathbf{D}_{\alpha}(q^{-1}_{\alpha\beta}hq_{\alpha\beta}).
\end{equation*}
因此
\begin{equation*}
\chi(h) = \sum_{\alpha}\left\{\sum_{\beta}\mathrm{Tr}\,\mathbf{D}_{\alpha}(q^{-1}_{\alpha\beta}hq_{\alpha\beta})\right\},
\tag{4.7.24}
\end{equation*}
其中对给定的 $\alpha$，对 $\beta$ 的求和限制在使 $q^{-1}_{\alpha\beta}hq_{\alpha\beta} \in \mathbf{K}_{\alpha}$ 的那些 $\beta$ 上。但 $q_{\alpha\beta} \in \mathbf{H}$，故 $q^{-1}_{\alpha\beta}hq_{\alpha\beta} \in \mathbf{L}_{\alpha}$（实际上它就是 $\mathbf{L}_{\alpha}$ 中使等式成立的元素），从而上式正是 $(\mathbf{D}_{\alpha} \downarrow \mathbf{L}_{\alpha}) \uparrow \mathbf{H}$ 的特征标。定理得证。

现在可以证明关于两个诱导表示直积的定理。

\noindent\textbf{定理 4.7.7}\quad 设 $\mathbf{D}$ 是 $\mathbf{K}$ 的表示，$C$ 是 $\mathbf{H}$ 的表示，则
\begin{equation*}
(\mathbf{D} \uparrow \mathbf{G}) \boxtimes (C \uparrow \mathbf{G}) \equiv \sum_{\alpha}\left\{\left(\mathbf{D}_{\alpha} \boxtimes C\right) \downarrow \mathbf{L}_{\alpha}\right\} \uparrow \mathbf{G}.
\tag{4.7.25}
\end{equation*}
由定理 4.7.3，
\begin{equation*}
(\mathbf{D} \uparrow \mathbf{G}) \boxtimes (C \uparrow \mathbf{G}) \equiv \left[\left\{(\mathbf{D} \uparrow \mathbf{G}) \downarrow \mathbf{H}\right\} \boxtimes C\right] \uparrow \mathbf{G},
\end{equation*}
由定理 4.7.6 它等价于
\begin{equation*}
\sum_{\alpha}\left\{\left[(\mathbf{D}_{\alpha} \downarrow \mathbf{L}_{\alpha}) \uparrow \mathbf{H}\right] \boxtimes C\right\} \uparrow \mathbf{G}.
\end{equation*}
现在 $C$ 是 $\mathbf{H}$ 的表示，$(\mathbf{D}_{\alpha} \downarrow \mathbf{L}_{\alpha}) \uparrow \mathbf{H}$ 是从 $\mathbf{L}_{\alpha}$ 诱导到 $\mathbf{H}$ 的表示，于是可再用定理 4.7.3，得到等价于
\begin{equation*}
\sum_{\alpha}\left\{\left[(C \downarrow \mathbf{L}_{\alpha}) \boxtimes (\mathbf{D}_{\alpha} \downarrow \mathbf{L}_{\alpha})\right] \uparrow \mathbf{H}\right\} \uparrow \mathbf{G}.
\end{equation*}
由于限制对内直积是可分配的、诱导是传递的（定理 4.7.4），最终得到
\begin{equation*}
\sum_{\alpha}\left\{\left(\mathbf{D}_{\alpha} \boxtimes C\right) \downarrow \mathbf{L}_{\alpha}\right\} \uparrow \mathbf{G},
\end{equation*}
定理得证。

我们把它写成 $\sum_{\alpha}E_{\alpha} \uparrow \mathbf{G}$ 的形式，其中 $E_{\alpha}$ 是 $\mathbf{L}_{\alpha} = \mathbf{H} \cap \mathbf{K}_{\alpha}$ 的表示。现在设 $\mathbf{M}$ 是 $\mathbf{G}$ 的子群，$B$ 是 $\mathbf{M}$ 的表示。若写 $\mathbf{G} = \sum_{\beta}\mathbf{M}m_{\beta}\mathbf{H}$ 与 $\mathbf{N}_{\alpha\beta} = \mathbf{M}_{\beta} \cap \mathbf{L}_{\alpha}$，则同定理 4.7.7 一样可证
\begin{equation*}
\begin{array}{l}
(B \uparrow \mathbf{G}) \boxtimes (E_{\alpha} \uparrow \mathbf{G}) \equiv \sum_{\beta}\left\{(B_{\beta} \boxtimes E_{\alpha}) \downarrow \mathbf{N}_{\alpha\beta}\right\} \uparrow \mathbf{G}\\
\equiv \sum_{\beta}\left\{(B_{\beta} \boxtimes D_{\alpha} \boxtimes C) \downarrow \mathbf{N}_{\alpha\beta}\right\} \uparrow \mathbf{G},
\end{array}
\end{equation*}
由此得到结论
\begin{equation*}
\begin{array}{l}
(B \uparrow \mathbf{G}) \boxtimes (\mathbf{D} \uparrow \mathbf{G}) \boxtimes (C \uparrow \mathbf{G})\\
\quad \equiv \sum_{\alpha}\sum_{\beta}\left\{(B_{\beta} \boxtimes D_{\alpha} \boxtimes C) \downarrow (\mathbf{M}_{\beta} \cap \mathbf{K}_{\alpha} \cap \mathbf{H})\right\} \uparrow \mathbf{G}.
\end{array}
\tag{4.7.26}
\end{equation*}
反复应用定理 4.7.7，式 (4.7.26) 显然可以推广，给出任意多个诱导表示直积的表达式。然而我们关心的是 $\mathbf{M} = \mathbf{K}^{i}$、$\mathbf{K} = \mathbf{K}^{j}$、$\mathbf{H} = \mathbf{K}^{l}$ 的特殊情形。

对空间群，$\mathbf{G}$ 就是整个空间群，$\mathbf{K}^{i}$、$\mathbf{K}^{j}$ 与 $\mathbf{K}^{l}$ 分别是与 $\mathbf{k}_{i}$、$\mathbf{k}_{j}$、$\mathbf{k}_{l}$ 相联系的小群 $\mathbf{G}^{\mathbf{k}_{i}}$、$\mathbf{G}^{\mathbf{k}_{j}}$、$\mathbf{G}^{\mathbf{k}_{l}}$。于是得到计算 $C^{ij, l}_{pq, r}$ 的公式
\begin{equation*}
\begin{array}{l}
C^{ij, l}_{pq, r} = \sum_{\alpha}\sum_{\beta}\dfrac{1}{|\mathbf{N}_{\alpha\beta}|}\sum_{\{\gamma \mid \mathbf{w}\} \in \mathbf{N}_{\alpha\beta}}\chi^{i}_{p}\left(\{\beta \mid \mathbf{v}\}^{-1}\{\gamma \mid \mathbf{w}\}\{\beta \mid \mathbf{v}\}\right)\\
\quad \chi^{j}_{q}\left(\{\alpha \mid \mathbf{u}\}^{-1}\{\gamma \mid \mathbf{w}\}\{\alpha \mid \mathbf{u}\}\right)\chi^{l*}_{r}\left(\{\gamma \mid \mathbf{w}\}\right),
\end{array}
\tag{4.7.27}
\end{equation*}
其中在式 (4.7.27) 中把 $\{\alpha \mid \mathbf{u}\}$ 认同为 $d_{\alpha}$。对平移 $\{E \mid \mathbf{t}\}$ 求和后显然可见：除非存在等价关系
\begin{equation*}
\beta\mathbf{k}_{i} + \alpha\mathbf{k}_{j} \equiv \mathbf{k}_{l},
\tag{4.7.28}
\end{equation*}
否则 $C^{ij, l}_{pq, r}$ 为零。把这个限制记为求和号上的撇，并把 $\{\gamma \mid \mathbf{w}\}$ 取为 $\mathbf{T}$ 在 $\mathbf{N}_{\alpha\beta}$ 中的典型陪集代表，得到
\begin{equation*}
\begin{array}{l}
C^{ij, l}_{pq, r} = \sum_{\alpha}'\sum_{\beta}'\dfrac{|\mathbf{T}|}{|\mathbf{N}_{\alpha\beta}|}\sum_{\{\gamma \mid \mathbf{w}\} \in \mathbf{N}_{\alpha\beta}/\mathbf{T}}\chi^{i}_{p}\left(\{\beta \mid \mathbf{v}\}^{-1}\{\gamma \mid \mathbf{w}\}\{\beta \mid \mathbf{v}\}\right)\\
\quad \chi^{j}_{q}\left(\{\alpha \mid \mathbf{u}\}^{-1}\{\gamma \mid \mathbf{w}\}\{\alpha \mid \mathbf{u}\}\right)\chi^{l*}_{r}\left(\{\gamma \mid \mathbf{w}\}\right).
\end{array}
\tag{4.7.29}
\end{equation*}
由于一般来说 $\mathbf{N}_{\alpha\beta}/\mathbf{T}$ 中的元素很少，且对 $\alpha$ 与 $\beta$ 的限制非常苛刻，三重求和后幸存的合适的 $\alpha$ 与 $\beta$ 很少多于一个，所以表达式 (4.7.29) 易于计算。

总结一下：要用式 (4.7.29) 计算 $C^{ij, l}_{pq, r}$，必须确定 $\mathbf{k}_{j}$ 关于 $\mathbf{G}^{\mathbf{k}_{i}}$ 的共星（co-star，即所有可能的 $\alpha\mathbf{k}_{j}$）、$\mathbf{k}_{i}$ 关于 $\mathbf{G}^{\mathbf{k}_{j}}$ 的共星（即所有可能的 $\beta\mathbf{k}_{i}$），以及交群 $\mathbf{N}_{\alpha\beta}$。

为了说明本节的结果，我们用空间群 $\mathbf{G} = P23\ (T^{1})$ 给出一个例子（Bradley 1966）。之所以选这个例子，是为了让式 (4.7.29) 中的双重求和出现不止一项幸存的情形。要找到对 $\alpha$ 与 $\beta$ 的求和非平凡的情形，看来必须考虑这样的空间群：其布里渊区的表示域（见定义 3.7.1）大于基本域（见定义 3.3.3），即其同形点群的阶低于相应晶系全对称点群的阶。空间群 $P23\ (T^{1})$ 建立在简立方布拉维格子上。该格子及其布里渊区的细节见第三章，算符所用记号与第一章定义的完全相同。$P23\ (T^{1})$ 的同形点群是四面体群 $23\ (T)$，它含 12 个元素 $E$、$C_{2m}$、$C^{\pm}_{3j}$（$m = x, y, z$；$j = 1, 2, 3, 4$）。我们只考虑布里渊区表示域的两个点（见图 3.13）：$\Gamma = (0, 0, 0)$ 与 $M = (\frac{1}{2}, \frac{1}{2}, 0)$。坐标照例以倒格子矢量 $\mathbf{g}_{1}$、$\mathbf{g}_{2}$、$\mathbf{g}_{3}$（沿 $k_{x}$、$k_{y}$、$k_{z}$ 方向，见表 3.3）为单位给出。定义点 $M^{+} = C^{+}_{31}M = (0, \frac{1}{2}, \frac{1}{2})$ 与 $M^{-} = C^{-}_{31}M = (\frac{1}{2}, 0, \frac{1}{2})$。$\Gamma$ 的星就是 $\Gamma$ 自己，而 $M$ 的星由三点 $M$、$M^{+}$、$M^{-}$ 组成。

\begin{table}[H]
\centering
\caption*{\textbf{表 4.4}\ $\mathbf{G}^{\Gamma}$ 与 $\mathbf{G}^{M}$ 的小表示}
\begin{tabular}{@{}lcccc@{}}
\toprule
$\mathbf{G}^{\Gamma}$ & $E$ & $3C_{2m}$ & $4C^{+}_{3j}$ & $4C^{-}_{3j}$\\
\midrule
$A$ & 1 & 1 & 1 & 1\\
${}^{1}E$ & 1 & 1 & $\omega^{*}$ & $\omega$\\
${}^{2}E$ & 1 & 1 & $\omega$ & $\omega^{*}$\\
$T$ & 3 & $-1$ & 0 & 0\\
\bottomrule
\end{tabular}
\quad
\begin{tabular}{@{}lcccc@{}}
\toprule
$\mathbf{G}^{M}$ & $E$ & $C_{2x}$ & $C_{2y}$ & $C_{2z}$\\
\midrule
$A_{1}$ & 1 & 1 & 1 & 1\\
$B_{1}$ & 1 & $-1$ & $-1$ & 1\\
$B_{2}$ & 1 & $-1$ & 1 & $-1$\\
$B_{3}$ & 1 & 1 & $-1$ & $-1$\\
\bottomrule
\end{tabular}
\end{table}

\noindent\textbf{表 4.4 的注}\quad（i）$\omega = \exp(2\pi i/3)$。（ii）$\mathbf{G}^{\Gamma}$ 中所有平移都由恒等元表示。（iii）在 $\mathbf{G}^{M}$ 中，$\{E \mid \mathbf{t}_{1}\}$ 与 $\{E \mid \mathbf{t}_{2}\}$ 由 $-1$ 表示，$\{E \mid \mathbf{t}_{3}\}$ 由 $+1$ 表示。

在表 4.4 中我们列出了小群 $\mathbf{G}^{\Gamma}$ 与 $\mathbf{G}^{M}$ 的小表示的特征标。$\Gamma$ 的小群是整个空间群，而 $M$ 的小群只含与四个转动算符 $E$、$C_{2x}$、$C_{2y}$、$C_{2z}$ 相结合的平移的倍数。

对我们的例子，考虑从 $\mathbf{G}^{M}$ 在 $\mathbf{G}$ 中诱导的表示对的内直积。我们用 $MA_{1}$ 记 $A_{1} \uparrow \mathbf{G}$，等等。按本节前文的记号，$\mathbf{k}_{j} = M = (\frac{1}{2}, \frac{1}{2}, 0)$ 且 $\mathbf{k}_{i} = M = (\frac{1}{2}, \frac{1}{2}, 0)$。很快可以验证：取这样的 $\mathbf{k}_{j}$ 与 $\mathbf{k}_{i}$，式 (4.7.1) 右端可能出现的 $\mathbf{k}_{l}$ 只有 $\mathbf{k}_{l} = \Gamma = (0, 0, 0)$ 与 $\mathbf{k}_{l} = M = (\frac{1}{2}, \frac{1}{2}, 0)$。

先考虑 $\mathbf{k}_{l} = \Gamma$ 的情形。由式 (4.7.19)，取 $\mathbf{H} = \mathbf{G}^{\Gamma}$、$\mathbf{K} = \mathbf{G}^{M}$，很快可以验证第一个双陪集分解中只有一项，唯一的代表 $d_{\alpha} = \{\alpha \mid \mathbf{0}\}$ 可取为恒等元 $\{E \mid \mathbf{0}\}$。于是 $\mathbf{K}_{\alpha} = \mathbf{G}^{M}$，交群 $\mathbf{L}_{\alpha} = \mathbf{G}^{\Gamma} \cap \mathbf{G}^{M} = \mathbf{G}^{M}$。接着作第二个双陪集分解 $\mathbf{G} = \sum_{\beta}\mathbf{G}^{M}b_{\beta}\mathbf{G}^{M}$，这次可以验证有三项幸存，三个代表 $b_{\beta} = \{\beta \mid \mathbf{0}\}$ 可取为 $\{E \mid \mathbf{0}\}$、$\{C^{+}_{31} \mid \mathbf{0}\}$ 与 $\{C^{-}_{31} \mid \mathbf{0}\}$。于是集合 $(\alpha\mathbf{k}_{j})$ 只由点 $M$ 组成，集合 $(\beta\mathbf{k}_{i})$ 由三点 $M$、$M^{+}$、$M^{-}$ 组成。经受住式 (4.7.28) 限制的唯一配对是 $\alpha = \beta = E$（即 $M + M \sim \Gamma$）。三重交群 $\mathbf{N}_{\alpha\beta} = \mathbf{G}^{M} \cap \mathbf{G}^{M} = \mathbf{G}^{M}$。

再考虑 $\mathbf{k}_{l} = M$ 的情形。这次第一个双陪集分解含三项，相应的 $\alpha$ 是 $E$、$C^{+}_{31}$ 与 $C^{-}_{31}$。由于 $\mathbf{G}^{M}$ 在这三个操作下不变，群 $\mathbf{G}^{M}$、$\mathbf{G}^{M+}$ 与 $\mathbf{G}^{M-}$ 相同，故 $\mathbf{L}_{\alpha} = \mathbf{G}^{M}$ 对三个 $\alpha$ 都成立。第二个双陪集分解与上一段相同（三项），而这次经受住式 (4.7.28)（取 $\mathbf{k}_{l} = M$）的配对不只是一个。例如配对 $\alpha\mathbf{k}_{j} = M^{-}$、$\beta\mathbf{k}_{i} = M^{+}$、$\mathbf{k}_{l} = M$ 满足 $M^{-} + M^{+} \equiv M$。此时式 (4.7.29) 的相应部分是 $\frac{1}{4}\left\{\chi^{M}_{r}(E) - \chi^{M}_{r}(C_{2x}) - \chi^{M}_{r}(C_{2y}) + \chi^{M}_{r}(C_{2z})\right\}$，当 $r = A_{1}$、$B_{2}$ 或 $B_{3}$ 时为零，当 $r = B_{1}$ 时为 1。汇总这些结果，与式 (4.7.1) 相应的方程是
\begin{equation*}
MB_{2} \boxtimes MB_{3} \equiv \Gamma T + MA_{1} + MB_{1}.
\tag{4.7.30}
\end{equation*}
作为式 (4.7.30) 的检验，注意两端维数都是 9（记住从 $\mathbf{G}^{M}$ 在 $P23$ 中诱导的表示维数为 3）。在表 4.5 中我们列出从 $\mathbf{G}^{M}$ 诱导的表示的全部 10 个这样的乘积。

\begin{table}[H]
\centering
\caption*{\textbf{表 4.5}\ $P23$ 中属于 $M$ 的表示的内直积}
\begin{tabular}{@{}ll@{}}
\toprule
$MA_{1} \boxtimes MA_{1} \equiv \Gamma A + \Gamma{}^{1}E + \Gamma{}^{2}E + 2MA_{1}$ & $MB_{1} \boxtimes MB_{2} \equiv \Gamma T + MB_{3} + MA_{1}$\\
$MA_{1} \boxtimes MB_{1} \equiv \Gamma T + MB_{2} + MB_{3}$ & $MB_{1} \boxtimes MB_{3} \equiv \Gamma T + MB_{2} + MA_{1}$\\
$MA_{1} \boxtimes MB_{2} \equiv \Gamma T + MB_{3} + MB_{1}$ & $MB_{2} \boxtimes MB_{2} \equiv \Gamma A + \Gamma{}^{1}E + \Gamma{}^{2}E + 2MB_{2}$\\
$MA_{1} \boxtimes MB_{3} \equiv \Gamma T + MB_{1} + MB_{2}$ & $MB_{2} \boxtimes MB_{3} \equiv \Gamma T + MB_{1} + MA_{1}$\\
$MB_{1} \boxtimes MB_{1} \equiv \Gamma A + \Gamma{}^{1}E + \Gamma{}^{2}E + 2MB_{1}$ & $MB_{3} \boxtimes MB_{3} \equiv \Gamma A + \Gamma{}^{1}E + \Gamma{}^{2}E + 2MB_{3}$\\
\bottomrule
\end{tabular}
\end{table}

\noindent\textbf{表 4.5 的注}\quad（i）表的每一行都是形如式 (4.7.30) 的一个方程。

（ii）为得到 Clebsch--Gordan 系数 $C^{ij, l}_{pq, r}$：选出左端含有符号 $p$ 与 $q$ 的那个方程，取其右端与 $r$ 对应的符号前面的数字。例如 $C^{MM, M}_{A_{1}A_{1}, A_{1}} = 2$。

\subsection*{晶体中的空间群选择定则}

本节所述理论在物理术语上的重要性，在于确定离子晶体、金属或半导体固体中各种量子力学过程的选择定则。这类过程的一个例子是：固体样品吸收一个光子并在固体中产生两个声子（波矢为 $\mathbf{k}$ 与 $-\mathbf{k}$ 以保持动量守恒）。其他过程可能涉及吸收一个光子并产生两个磁振子（一个在 $\mathbf{k}$、一个在 $-\mathbf{k}$），或产生一个声子（在 $\mathbf{k}$）与一个磁振子（在 $-\mathbf{k}$），或只产生一个声子（$\mathbf{k} = 0$）或一个磁振子（$\mathbf{k} = 0$）。在研究固体中量子力学粒子或准粒子之间的相互作用时（例如电子—声子相互作用或磁振子—声子相互作用）也会出现选择定则。确定某个过程是否被允许的选择定则，涉及空间群表示的直积分解（或可能涉及空间群表示的对称化直积分解，见 4.8 节）。设 $\Psi^{i}_{r}$ 是属于表示 $\Gamma^{i}$ 第 $r$ 行的波函数，$\mathbf{W}$ 是属于表示 $\Gamma^{W}$ 的自伴算符。初态 $\Psi^{i}_{r}$ 与末态之间的矩阵元一般非零，当且仅当直积 $(\Gamma^{i} \uparrow \mathbf{G}) \boxtimes \Gamma^{W} \boxtimes (\Gamma^{f*} \uparrow \mathbf{G})$ 中含全对称表示；这正是本节理论所能判断的。
